% Source: http://tex.stackexchange.com/a/150903/23931
\documentclass{article}
\usepackage[letterpaper,margin=1in]{geometry}
\usepackage{xcolor}
\usepackage{fancyhdr}
\usepackage{tgschola} % or any other font package you like

\pagestyle{fancy}
\fancyhf{}
\fancyhead[C]{%
  \footnotesize\sffamily
  \yourname\quad
  web: \textcolor{blue}{\itshape\yourweb}\quad
  \textcolor{blue}{\youremail}}

\newcommand{\soptitle}{Statement of Purpose}
\newcommand{\yourname}{Hrishikesh Belagali}
\newcommand{\youremail}{belagal1@msu.edu}
\newcommand{\yourweb}{https://lonelyneutrin0.github.io/portfolio/}

\newcommand{\statement}[1]{\par\medskip
  \underline{\textcolor{blue}{\textbf{#1:}}}\space
}
\vspace*{-20pt}
\usepackage[
  colorlinks,
  breaklinks,
  pdftitle={\yourname - \soptitle},
  pdfauthor={\yourname},
  unicode
]{hyperref}
\parindent=0pt
\begin{document}

\begin{center}\LARGE\soptitle\\
\large of \yourname\
\end{center}

\hrule
\vspace{1pt}
\hrule height 1pt

\bigskip
% \statement{Main Statement}
% \textbf{I am driven by the challenge of modelling complex physical systems. I started doing this through molecular dynamics, where I learnt
% about the power of computational modelling, as well as some of the limitations of classical resources. I then became interested in quantum algorithms
% and the new possibilities they open up for modelling. Alongside formal research, I explored various computational modelling and optimization methods 
% such as classical Monte Carlo methods, quantum optimization, and machine learning prediction for dynamical systems. I aim to 
% go to graduate school in quantum algorithms and contribute to research in developing novel methods for simulation and optimization.} \\\\


% \statement{Introduction}
% \begin{itemize}
%   \item Inrtoduce your fascination with modelling complex physical systems.
%   \item Emphasize that this curiosity has driven your academic and research pursuits. 
%   \item Set the state for classical $\to$ limitations $\to$ quantum.
% \end{itemize}
% \statement{Classical Simulation: Power and limits}
% \begin{itemize}
%   \item Early exposure to molecular dynamics simulations in NAMD and VMD
%   \item Systems studied: Thylakoid membrane protein interactions, aquaporin capping mechanisms, small molecule effect on membrane dynamics and stability
%   \item Work done: setting up systems, running on clusters, analyzing results 
%   \item Taught about the power of computational modelling, and also limitations of classical resource (timescales, system size)
% \end{itemize}
% \statement{Why Quantum Algorithms}
% \begin{itemize}
%   \item Interest in quantum computing sparked by the limitations of classical simulations
%   \item A different perspective on the same problem with novel insights 
%   \item MSU-Q Honors Research Scholar program - generalized quantum simulators, compilation optimization, group-theoretic methods
%   \item Mention programming tools (Python, Cirq)
%   \item Mention QuTiP contribution - understanding how quantum systems are modelled in Python packages
% \end{itemize}
% \statement{Broader Computational Physics Projects}
% \begin{itemize}
%   \item Take Monte Carlo, Quantum Optimization, machine learning projects.
%   \item Different ways of modelling and optimizing complex systems 
%   \item Talk about the skills you gained here
% \end{itemize}
% \statement{Conclusion}
% All experiences point towards modelling complex physical systems using various computational methods, out of which 
% quantum algorithms is the most exciting and promising. I want to pursue graduate studies in this field, and contribute to research 
% in developing novel quantum algorithms for simulation and optimization, which has significant applications in materials science and 
% pharmaceuticals.
% \newpage
My research interests are centered around a single question: how can we model complex physical systems effectively? 
I always thought science was about understanding the natural world through observation and experimentation. 
My research experience has revealed that computation is an equally powerful tool, which can provide unique insights inaccessible to traditional experiments.
My interest in answering this question has guided me across multiple computational frameworks, including molecular simulation, 
quantum algorithms, and machine learning techniques. By participating in the Science Undergraduate Laboratory Internship, I aim to deepen my understanding of 
quantum simulation and optimization methods by learning from experts in the field. This program will help me explore the field of 
quantum algorithms for molecular simulation and understand how such methods can address computational bottlenecks, 
which will be invaluable as I consider PhD programs in quantum computing.
\\

As an undergraduate research assistant in the Vermaas Lab, I work on molecular dynamics simulations of biomolecular systems 
using VMD and NAMD. Over the past year, I have been working on characterizing water transport mechanisms in aquaporin channels in plants, which are vital for maintaining cellular water balance 
under drought stress. Understanding this mechanism is essential for bioengineering efforts to improve plant drought 
tolerance. During this study, I faced challenges quantifying the transport mechanism because of the large system size, 
which restricted the simulation length and sampling. I explored alternative methods 
for simulation such as grid-steered molecular dynamics, which taught me to consider computational constraints while studying 
complicated systems.
\\

My experience with the aquaporin project highlighted the limitations of classical simulation. Simulating larger systems, capturing statistically 
rare events, and reaching biologically significant timescales is computationally expensive, which limits the scope of 
phenomena we can study. These limitations led me to explore alternative computational approaches, of which quantum computing stood out to me 
because it leverages principles of quantum mechanics to solve complex problems. In order to explore my interest in 
quantum computing, I work as an Honors Research Scholar under Professor Ryan LaRose with a focus on generalized quantum 
simulators. My research involves studying group-theoretic methods and compilation optimization techniques in Python and Cirq to develop
methods for simulating quantum algorithms on classical computers. The development of advanced quantum simulation techniques is useful 
for benchmarking and refining larger-scale quantum algorithms, which might not be feasible to run on current quantum hardware. By 
contributing to this research, I've learned to consider both theoretical performance and practical challenges of 
implementing quantum algorithms.
\\

Alongside formal research, I have built several computational physics projects using Python, NumPy, Matplotlib, and Qiskit.
In my first semester of college, I was exposed to Monte Carlo algorithms, and was amazed by how directed random search 
could be used to solve a variety of difficult problems. I implemented a simulated annealing algorithm for 
Quadratic Unconstrained Binary Optimization (QUBO) problems and later wrote a quantum optimization algorithm using adiabatic evolution in Qiskit. 
Experimenting with different evolution schedules and total evolution times helped me understand the limitations of 
current quantum simulators and improved my scientific programming skills. 
\\

I am eager to further explore the intersection of molecular simulation and quantum algorithms through the SULI program at Oak Ridge 
National Laboratory. ORNL's focus on biophysics and simulations, the Quantum Science Center's QHPC initiative and research on hybrid quantum-classical 
algorithms closely align with my interests. I am particularly excited about the prospect of working under the guidance of 
Dr. Ryan Bennink, whose research on quantum high performance computing is very intriguing.
ORNL's interdisciplinary research environment will allow me to improve my scientific communication and collaboration skills,
which will be beneficial for my career. I believe that my intersection of molecular dynamics and quantum algorithms research experience 
will allow me to meaningfully contribute to ongoing projects at ORNL, and the skills that I gain through this program will prepare me for graduate school and a career in research.
In the long term, I'm interested in working on methods that bridge molecular simulation and quantum algorithms to model materials 
and biomolecular systems effectively.
\end{document}
